\chapter{CONCLUSION}

\section{Conclusion}

This project successfully developed MorphAI, an identity-preserving facial
inpainting system based on the LGNet architecture. The proposed system
reconstructs masked facial regions while maintaining facial structure, local
textures, and the identity of the input subject. The pipeline combines facial
landmark extraction, region-specific mask generation, landmark heatmaps, and
a five-channel input representation for the LGNet generator.

\noindent The model was trained for five epochs using mixed-precision training on a
dataset containing more than 16,800 facial images. The training loss
decreased consistently, with the fifth epoch selected as the final
convergence checkpoint. Evaluation was performed on 16,800 test samples
across different facial regions and compared with the LaMa baseline.

\noindent The proposed LGNet achieved an overall PSNR of 40.99 dB, SSIM of 0.9928,
ArcFace identity score of 0.9819, RMSE$_{log}$ of 0.0069, and FID of 0.1398.
Compared with LaMa, LGNet provided better structural and identity
preservation, particularly for challenging regions such as the mouth and
nose. These results demonstrate that the proposed approach can effectively
restore facial regions while maintaining the visual characteristics and
identity of the original face.

\section{Limitations}

Despite the promising results, the system has several limitations:

\begin{itemize}
    \item The system is primarily designed for 2D facial image inpainting
    and does not currently reconstruct or represent faces in 3D.
    
    \item The current implementation is not designed as a live or
    real-time video inpainting system and is mainly intended for
    single-image reconstruction.
    
    \item The quality of reconstruction depends on the accuracy of facial
    landmark detection and the generated region mask.
    
    \item Performance may decrease for heavily corrupted regions or
    facial structures with limited surrounding visual information.
    
    \item The model requires GPU resources for efficient inference and
    training, which may limit deployment on low-resource devices.
    
    \item Although ArcFace similarity is used to evaluate identity
    preservation, the identity score does not guarantee perfect
    preservation in every reconstructed image.
\end{itemize}

\section{Future Works}

Based on the findings of this project, several improvements can be explored
in future work:

\begin{itemize}
    \item \textbf{3D Facial Reconstruction:} Extend MorphAI from 2D image
    inpainting to 3D facial reconstruction to model facial geometry and
    provide more realistic results from different viewpoints.
    
    \item \textbf{Real-Time Video Inpainting:} Extend the system to process
    video streams and reconstruct masked facial regions continuously while
    maintaining temporal consistency between frames.
    
    \item \textbf{Dataset Expansion:} Train the model on larger and more
    diverse facial datasets containing different poses, lighting
    conditions, expressions, and degrees of facial occlusion.
    
    \item \textbf{Improved Identity Preservation:} Explore stronger
    identity-aware objectives and modern face-recognition embeddings to
    further improve preservation of facial identity.
    
    \item \textbf{Advanced Generative Models:} Investigate diffusion-based
    and other modern generative architectures to improve reconstruction
    quality for complex facial regions.
    
    \item \textbf{Mobile and Edge Deployment:} Optimize the model through
    techniques such as quantization, pruning, or knowledge distillation
    to enable efficient deployment on resource-constrained devices.
    
    \item \textbf{Interactive Reconstruction:} Develop a more flexible
    interface that allows users to select, adjust, and preview different
    facial regions before generating the final reconstruction.
\end{itemize}